% =====================================================================
%  CONSIDERAÇÕES FINAIS
%
%  Revisado em 2026-10-05. A estrutura da FATEC-ID é Introdução, Capítulos
%  I a III e Considerações Finais, então esta parte não é "Capítulo IV":
%  fica sem número, como a introdução, e em texto corrido.
%  O Manual (item 4.1.2) exige que ela (a) responda à pergunta-problema,
%  (b) revisite cada hipótese confirmando ou refutando e (c) diga se os
%  objetivos foram alcançados. Os parágrafos seguem essa ordem.
%  Todos os números já aparecem no Capítulo III, com a procedência anotada lá.
%  Revisto em 2026-10-06 depois da refação dos experimentos (EA corrigido).
%  Versão anterior em notas/rascunhos/antes_revisao_2026-10-05/.
% =====================================================================
\chapter*{Considerações finais}
\phantomsection
\addcontentsline{toc}{chapter}{CONSIDERAÇÕES FINAIS}

Este trabalho investigou o efeito de acrescentar um modelo de volatilidade
condicional a uma estratégia automatizada de negociação cuja previsão de direção
permanece a mesma. A questão norteadora foi em que medida a integração dos modelos
ARIMA e GARCH em um único algoritmo de negociação altera o desempenho e o risco
das operações, em comparação com o uso do ARIMA isolado e com uma estratégia de
cruzamento de médias móveis.

Quanto ao risco, a resposta é clara. Em todos os períodos analisados, a versão do
algoritmo com o GARCH sofreu quedas menores que a sua gêmea, que usa o mesmo ARIMA
sem o modelo de volatilidade, e o efeito se repetiu nas outras três ações
estudadas, sobretudo nos períodos de turbulência. O ano de 2008 resume bem esse
comportamento. Enquanto a PETR4 perdia quase metade do seu valor, o ARIMA-GARCH
encerrou o ano com um pequeno ganho.

Quanto ao desempenho, o efeito foi menos claro. Na maior parte dos testes, a
versão com GARCH terminou com retorno maior que o da gêmea, mas essa vantagem só
foi confirmada estatisticamente em uma das comparações, e em nenhuma delas o
GARCH piorou o resultado de forma significativa. Em relação às médias
móveis de \citeonline{almeida2016}, o ARIMA-GARCH protegeu melhor nos momentos de
queda, mas ganhou em alguns períodos e perdeu em outros, por diferenças que ficaram
dentro da margem do acaso. Em relação à taxa Selic, o ARIMA-GARCH só rendeu mais
no período de 2017 a 2024, e por uma margem que o acaso também explicaria. Nos
demais períodos, rendeu menos.

A primeira hipótese afirmava que a integração entre ARIMA e GARCH gera sinais
mais eficazes que o ARIMA isolado, tanto em acerto direcional quanto em adequação
ao risco. Ela foi parcialmente confirmada. A adequação ao risco foi confirmada
pelas quedas menores em todos os períodos, ainda que o desempenho ajustado ao
risco, medido em relação à Selic, só tenha melhorado de forma significativa em
uma das comparações. O acerto direcional, por outro lado, não
melhorou e permaneceu próximo do acaso. Esse resultado é coerente com a teoria,
pois o GARCH descreve o tamanho das oscilações, e não a direção dos preços. O ganho
veio da seletividade. Ao deixar de operar nos dias de maior volatilidade prevista,
o algoritmo agiu menos vezes e em condições mais regulares.

A segunda hipótese afirmava que a combinação representa um avanço em relação às
análises baseadas em médias móveis. Ela foi confirmada apenas quanto ao risco. Em
todos os períodos, houve uma variante com GARCH que protegeu mais do que as médias
móveis, e a diferença foi grande sobretudo nos períodos de crise. Ainda assim,
nem toda variante protegeu mais, e as que mantêm a posição aberta durante o veto
chegaram a ter quedas maiores. Quanto ao retorno, as diferenças entre os dois
algoritmos ficaram dentro da margem do acaso.

O objetivo geral, desenvolver em MQL5 um \textit{Expert Advisor} que integra os
modelos ARIMA e GARCH e avaliar experimentalmente o efeito da inclusão do GARCH,
foi alcançado. Os dois modelos foram implementados na própria linguagem da
plataforma e validados por recuperação de parâmetros em séries com resposta
conhecida. A escolha das ordens do ARIMA foi automatizada pelo teste ADF e pelo
critério de Akaike. A estratégia de \citeonline{almeida2016} foi reimplementada,
e a reprodução dos seus testes confirmou o comportamento do algoritmo original,
embora não os seus valores em reais, que dependem de informações que o estudo
original não registra. Por fim, o protocolo com períodos de calibração e de teste
separados foi definido antes da observação dos resultados e executado sobre dados
oficiais, e os três algoritmos foram comparados quanto a retorno, risco e
desempenho ajustado ao risco, com testes estatísticos.

Os resultados têm as limitações discutidas ao final do Capítulo III. A maior parte
das diferenças não atinge significância estatística, a configuração que se sai
melhor em um período não tende a repetir o desempenho no seguinte, os ganhos se
concentram em poucos anos, a análise se restringe a quatro ativos de um único
mercado e a medida de desempenho não inclui os proventos, que, quando incluídos,
favorecem as estratégias que permanecem mais tempo compradas. Além disso, os experimentos foram realizados por simulação sobre dados
históricos, com execução no preço de abertura e sem considerar a falta de liquidez
ou a diferença entre o preço pedido e o executado, condições que só podem ser
verificadas em operação real.

Como continuação, o resultado indica que o caminho mais promissor está no
controle de risco, que foi onde o modelo mostrou efeito. Uma possibilidade é
empregar variações assimétricas do GARCH, que reagem mais às quedas do que às
altas, comportamento comum nas séries de ações e não capturado pela especificação
simétrica utilizada.
% REFERÊNCIA: Glosten, Jagannathan e Runkle (1993), Nelson (1991)
Outra é aplicar o mesmo protocolo a um conjunto maior de ativos e acompanhar o
algoritmo em conta de demonstração, em tempo real, o que permitiria também avaliar
as ordens de limite de perda e de realização de lucro que
\citeonline{almeida2016} previu em seu algoritmo, mas não utilizou em seus testes.

No eixo da previsão, a próxima etapa é a que estava prevista no projeto original
deste trabalho e foi deixada de lado para que a implementação do ARIMA e do GARCH
pudesse ser concluída e validada: substituir o ARIMA por uma rede neural
recorrente do tipo LSTM, do inglês \textit{Long Short-Term Memory}, capaz de representar relações não lineares, mantendo o
GARCH como filtro de risco. A arquitetura do algoritmo, que separa a previsão da
direção da avaliação do risco, permite essa troca sem alterar o restante. Os
resultados deste trabalho, porém, recomendam cautela na expectativa, já que a
taxa de acerto direcional do ARIMA ficou próxima do acaso e o ganho observado veio
do controle de risco, e não da previsão.
